Standard large language model pretraining samples uniformly across all available data, ignoring the potential benefits of strategic data ordering. We investigate whether introducing controlled distribution shifts during pretraining---by withholding specific domains and introducing them later---produces models that adapt more rapidly to new data distributions. Using a 30M-parameter GPT-2 model trained on 13 domains from The Pile, we compare a two-phase \emph{distribution shift curriculum} (training first on 9 core domains, then expanding to include 3 withheld shift domains) against a standard uniform baseline. We find three key results: (1) the curriculum-trained model adapts \textbf{2.8$\times$ faster} to newly introduced domains during a 500-step adaptation phase (9.1 vs.\ 3.3 perplexity improvement); (2) the curriculum model \emph{unexpectedly outperforms} the baseline on 7 of 9 core domains it trained on throughout, with 5 reaching statistical significance; and (3) the curriculum model shows 26\% less catastrophic forgetting on held-out domains during adaptation. However, the curriculum model finishes training with higher perplexity on shift domains and shows no advantage on completely unseen domains. These results provide partial but meaningful support for the distribution shift curriculum hypothesis, with the strongest evidence for improved adaptation speed---a property with practical implications for deploying language models in specialized domains.
